\documentclass[11pt]{article}
\usepackage[a4paper,margin=1in]{geometry}
\usepackage{amsmath,amssymb,amsthm,mathtools,hyperref}
\setlength{\emergencystretch}{1em}
\newtheorem{theorem}{Theorem}
\newtheorem{lemma}[theorem]{Lemma}
\newcommand{\dd}{\,\mathrm d}
\newcommand{\norm}[1]{\left\lVert#1\right\rVert}
\title{Global null control of a two-phase Stefan interface with signed initial temperatures}
\author{Kaiwen Jin\thanks{\texttt{yui.ui.siki@gmail.com}}}
\date{October 12, 2026}
\begin{document}
\maketitle

\begin{abstract}
We prove that two Dirichlet boundary controls can drive a one-dimensional
two-phase Stefan system to zero temperature and any prescribed interior
interface position in every positive time. The initial temperatures may be
arbitrary real smooth functions compactly supported in their phases, and the
interface remains strictly inside the rod throughout. A short uncontrolled
evolution gives analytic interface traces without a sign assumption. These
traces are then extended by Gevrey interpolation and the flatness series of
Colle, Loh\'eac and Takahashi, while a convex interpolation of the interface
position prevents collision with the endpoints.
\end{abstract}

\section{The control problem}
In the two-phase Stefan system, heat diffusion acts on both sides of a
moving interface, and the difference of the interface heat fluxes determines
its velocity. Boundary-temperature control asks whether the two exterior
temperatures can be chosen to bring both phases to zero temperature and
place the interface at a prescribed point. The geometric constraint is
essential: neither phase may disappear during the control.

Colle, Loh\'eac and Takahashi \cite[Theorem 1.1]{CLT} prove arbitrary-time
boundary controllability by a flatness construction for initially
phase-signed temperatures, with strict signs at the exterior endpoints.
The coefficient estimates in their Proposition 3.4 and Corollary 3.6
do not require these signs or a smallness condition. The signs enter their
preparatory regularization argument. We instead prove short-time
analyticity directly for arbitrary signed compactly supported profiles.
This permits the same series reconstruction to start from the actual
positive-time solution.

The target here is the unconstrained-sign formulation in AIM Problem 258
\cite{AIM258}. The temperature fields are treated as real solutions of the
specified heat and Stefan equations; no phase-sign inequality is added.
Let $L$ denote the rod length, $d_l,d_r$ the diffusivities and $\lambda$
the latent-heat coefficient. All are positive. Let $L_0,L_T$ be the initial
and desired interface positions. Time is the first argument of every
temperature below.

\begin{theorem}\label{thm:control}
For every $L,d_l,d_r,\lambda,T>0$, every $L_0,L_T\in(0,L)$,
and every real $y_0\in C_c^\infty(0,L_0)$ and
$z_0\in C_c^\infty(L_0,L)$, there exist $v_l,v_r\in L^\infty(0,T)$,
$\ell\in C^1([0,T])$ with $0<\ell<L$, and temperatures satisfying
\[
 y_t=d_l y_{xx}\quad(0<x<\ell(t)),\qquad
 z_t=d_r z_{xx}\quad(\ell(t)<x<L),
\]
\[
 y(t,0)=v_l(t),\quad z(t,L)=v_r(t),\quad
 y(t,\ell(t))=z(t,\ell(t))=0,
\]
\[
 d_l y_x(t,\ell(t)^-)-d_r z_x(t,\ell(t)^+)=-\lambda\ell'(t).
\]
The initial data and interface are $(y_0,z_0,L_0)$, the terminal
temperatures vanish, and $\ell(T)=L_T$. Zero extensions of both
temperatures must belong to $C([0,T];L^2(0,L))$. No temperature sign,
smallness condition, or uniform control-cost bound is imposed.
\end{theorem}

The construction gives a classical solution for positive time, with continuous
temperatures and spatial first derivatives in phase coordinates down to zero.
All constants may depend on the entire parameter and initial-data tuple.


The result allows arbitrarily large data and arbitrarily small control
times; it supplies no common bound on the control amplitudes. The controlled
temperatures need not retain phase signs. The initial class remains smooth
and compactly supported, and no higher-dimensional statement is made.

The proof first solves the system for a short time on two fixed intervals.
Complex-time heat estimates give an analytic interface and analytic flux
traces at positive time. The PDE then supplies spatial analyticity, so the
actual solution agrees with its interface-centered flatness series. We
keep this agreement on an open time interval and cut off the traces
afterwards. The interface is a convex combination of its short-time path
and the target position, which keeps both phases nonempty.

The following sections give the normalization, the short-time construction,
the series reconstruction, and the final control. Complex scalars occur only
in the analyticity argument; the final solution and controls are real.

\section{Normalization}
Put $r=x/L$, $b=\ell/L$, $b_0=L_0/L$, $b_T=L_T/L$, and
\[
 c_-=d_l/L^2,\qquad c_+=d_r/L^2,
 \qquad \theta_-(t,r)=-\frac{d_l}{\lambda L^2}y(t,Lr),\qquad
 \theta_+(t,r)=-\frac{d_r}{\lambda L^2}z(t,Lr).
\]
Then $\theta_{\pm,t}=c_\pm\theta_{\pm,rr}$ in their phases,
the interface values vanish, and
\[
 b'=\theta_{-,r}(t,b^-)-\theta_{+,r}(t,b^+). \tag{1}
\]
Indeed the right side is
$[-d_l y_x+d_r z_x]/(\lambda L)=\ell'/L$.
Conversely these identities recover exactly the physical equations and
the controls by multiplying each phase by $-\lambda L^2/d_\pm$.
It therefore suffices to work on $[0,1]$ with arbitrary $c_\pm>0$.

\section{Short-time solutions and analyticity}
\begin{lemma}[Spectral estimates]\label{lem:spectral}
Let $e_n(s)=\sqrt2\sin(n\pi s)$ and $\mu_n=(n\pi)^2$.
For $0\le\gamma$, set
\[
 E_\gamma=\{u=(u_n):\ \norm{u}_{E_\gamma}^2
             =\sum_{n\ge1}\mu_n^{2\gamma}|u_n|^2<\infty\}.
\]
These are complex Banach spaces. Identify $E_0$ with the completion of
the sine polynomials under the integral $L^2$ norm. Every continuous
function on $[0,1]$ is represented in this completion by its sine
coefficients. For $\alpha=7/8$, $E_\alpha$ embeds continuously into
$C^1([0,1])$, with zero endpoint values. For $\Re z>0$, define
\[
 S(z)u=(e^{-\mu_n z}u_n)_{n\ge1},\qquad
 Au=(\mu_nu_n)_{n\ge1},\quad D(A)=E_1.
\]
If $\Re z\ge k|z|$ for a fixed $k>0$, then
\[
 \norm{S(z)}_{E_\alpha\to E_\alpha}\le1,\quad
 \norm{S(z)}_{E_0\to E_\alpha}\le C_k|z|^{-\alpha},\quad
 \norm{S'(z)}_{E_0\to E_\alpha}\le C_k|z|^{-\alpha-1}. \tag{2}
\]
Moreover $S(z)$ is holomorphic away from zero, strongly continuous
at zero within such a sector, and if $u\in E_{\alpha+1}$,
$\norm{(S(z)-S(w))u}_{E_\alpha}\le |z-w|\norm{Au}_{E_\alpha}$
whenever the segment from $z$ to $w$ lies in $\Re z\ge0$.
\end{lemma}
\begin{proof}
Completeness of each weighted sequence space follows by taking coordinate
limits of a Cauchy sequence: finite partial sums pass to the limit, and
the Cauchy tail estimate bounds the full norm of each difference.
Direct integration gives $\int_0^1e_me_n\dd s=\delta_{mn}$.
For a continuous function, odd reflection about zero followed by
$2$-periodic extension is bounded and piecewise continuous, hence
Riemann integrable. Apply \cite[Theorem 11.8.13]{RA} after the change of variable $s\mapsto\pi s$:
the cosine coefficients vanish and the sine partial sums converge in
$L^2(0,1)$, with squared norm equal to the sum of squared coefficients.
This proves the asserted representation and its isometry on continuous
functions. No endpoint condition is required for this $L^2$ assertion.

For $u\in E_\alpha$, Cauchy--Schwarz for finite sums, followed by the
convergent numerical series $\sum n^{2-4\alpha}$, gives
\[
 \sum n|u_n|\le C\norm{u}_{E_\alpha}.
\]
Both the sine series and its derivative thus converge uniformly, proving
the embedding and endpoint statement. Inequalities (2) follow from
$\sup_{v\ge0}v^p e^{-kv}<\infty$ for $p=\alpha,\alpha+1$.
Finite spectral sums are holomorphic; their operator derivatives converge
locally uniformly away from zero by the same bounds. Strong continuity
at zero follows by splitting the sequence into a finite initial segment
and a uniformly small tail. Finally integrate
$S'(\zeta)u=-S(\zeta)Au$ along the indicated segment.
Compactly supported smooth functions belong to every $E_\gamma$,
by repeated integration by parts in their sine coefficients.

The scalar Cauchy formula and its Taylor, estimate and uniform-limit
consequences \cite[Theorems 3.2.15, 3.3.1, 3.3.8 and 3.4.5]{CA}
also hold for functions valued in these sequence spaces or
their finite products. The vector Cauchy formula follows by Riemann
integration on a circle and the scalar formula in each coordinate.
The integral's norm is bounded by length times the supremum norm.
Coordinate equality identifies the vectors. This proves the vector
Cauchy estimate and power expansion. The same formula proves that a
locally uniform vector limit of holomorphic functions is holomorphic.
An improper vector integral dominated in norm by an integrable scalar
function exists by the Cauchy criterion; if its truncated integrals are
holomorphic, the preceding limit argument applies.
\end{proof}

\begin{lemma}[Local existence and time analyticity]\label{lem:local}
For every $c_\pm>0$, $b_0\in(0,1)$ and real smooth compactly supported
initial profiles in the two phases, there is $\delta>0$ and a solution
with zero exterior controls on $[0,\delta]$. Its interface is $C^1$,
stays in $(0,1)$, and in fixed phase coordinates the profiles are
continuous into $E_\alpha$ down to zero and holomorphic in a complex
time sector at every positive real time. The interface and both interface
slopes are analytic in time on $(0,\delta)$.
\end{lemma}
\begin{proof}
Write
\[
 Y(t,s)=\theta_-(t,b(t)s),\qquad
 Z(t,s)=\theta_+(t,b(t)+(1-b(t))s).
\]
For $U=(Y,Z,b)$ define
\[
 F(U)=Y_s(1)/b-Z_s(0)/(1-b),\quad
 a_-(b)=c_-/b^2,\quad a_+(b)=c_+/(1-b)^2,
\]
\[
 G_-(U)=\frac{F(U)}b sY_s,\qquad
 G_+(U)=\frac{F(U)}{1-b}(1-s)Z_s.
\]
The transformed system is
\[
 b'=F(U),\quad Y_t=-a_-(b)AY+G_-(U),\quad
 Z_t=-a_+(b)AZ+G_+(U). \tag{3}
\]
The two signs of the drift follow by differentiating the two coordinate
maps at fixed $s$. Let $U_0=(Y_0,Z_0,b_0)$.
Lemma~\ref{lem:spectral} makes $F$ and $G_\pm$ holomorphic and locally Lipschitz, respectively
into $\mathbb C$ and $E_0$, on $E_\alpha^2\times\mathbb C$ where
$b\ne0,1$. Continuous spatial products use the isometry established in Lemma~\ref{lem:spectral}.

Choose $\vartheta=\pi/6$. On the closed time sector
$K_\delta=\{t:|t|\le\delta,\ |\arg t|\le\vartheta\}$, including its
tip, use the supremum product norm. Consider continuous functions,
holomorphic in the sector interior, with $U(0)=U_0$,
$\norm{Y-Y_0}_{E_\alpha},\norm{Z-Z_0}_{E_\alpha}\le1$ and
$|b-b_0|\le r$. Choose a positive $r$ sufficiently small that
$b,1-b$ are nonzero and
$|a_\pm(b)-a_\pm(b_0)|\le\varepsilon a_\pm(b_0)$ throughout this
complex disk, with a fixed sufficiently small $\varepsilon>0$.
This is a closed subset of a complete supremum space: completeness
and closure of holomorphy follow from Lemma~\ref{lem:spectral}. The rational functions and
their derivatives are uniformly bounded on this set.

Set $\tau_\pm(t)=\int_0^t a_\pm(b(q))\dd q$, using straight segments.
For $q=\rho t$, $0\le\rho<1$, one has
\[
 \tau_\pm(t)-\tau_\pm(q)
 =t\int_\rho^1a_\pm(b(ut))\dd u.
\]
Consequently these differences lie in a fixed closed sector of the
right half-plane and their modulus is comparable to $(1-\rho)|t|$.
The same is true of the straight segment joining these increments
for two inputs. Their difference is at most
$C(1-\rho)|t|\norm{b-\widehat b}_\infty$.

Define
\[
 \Phi_b(U)(t)=b_0+\int_0^tF(U(q))\dd q,
\]
\[
 \Phi_Y(U)(t)=S(\tau_-(t))Y_0+
 \int_0^t S(\tau_-(t)-\tau_-(q))G_-(U(q))\dd q, \tag{4}
\]
and the analogous expression for $\Phi_Z$.
Lemma~\ref{lem:spectral} gives a bound $C\delta^{1-\alpha}$ for each integral in (4).
It is continuous at the tip with value zero. Holomorphy follows by
writing it as $t\int_0^1\cdots\dd\rho$: truncations away from $\rho=1$
are holomorphic (Riemann sums and the vector Cauchy formula), and the
tail is locally uniformly bounded by $C\int(1-\rho)^{-\alpha}\dd\rho$.
The primitives starting at the tip can be justified by starting at
$\rho=\epsilon$ and letting $\epsilon\downarrow0$, with bounded tails.

For two inputs, integration of the $S'$ estimate in (2) on the straight
segment joining their increments proves
\[
 \norm{S(\tau(t)-\tau(\rho t))-
 S(\widehat\tau(t)-\widehat\tau(\rho t))}_{E_0\to E_\alpha}
 \le C\norm{b-\widehat b}_\infty[(1-\rho)|t|]^{-\alpha}.
\]
The initial terms differ by at most
$C\delta\norm{b-\widehat b}_\infty\norm{AY_0}_{E_\alpha}$,
and similarly for $Z_0$. Thus
\[
 \norm{\Phi(U)-\Phi(\widehat U)}_\infty
 \le C(\delta+\delta^{1-\alpha})\norm{U-\widehat U}_\infty. \tag{5}
\]
The same estimates against $U_0$ show that $\Phi$ maps the indicated
closed set into itself for sufficiently small $\delta$: the scalar
change is $O(\delta)$, and each profile change is
$O(\delta\norm{AY_0}_{E_\alpha}+\delta^{1-\alpha})$.
All constants are finite for these fixed, possibly large, data.
Choose $\delta$ also making (5) less than one. Iteration converges,
since successive differences are bounded by a summable geometric
series in a complete space. Passing to the limit in (4) gives a fixed
point; (5) gives uniqueness within this set. Conjugation preserves the
map and the initial data, so the fixed point is real on the real axis.

At positive real time, its $E_\alpha$ derivatives exist by holomorphy.
Each sine coefficient in (4) solves
$Y_n'=-\mu_n a_-Y_n+(G_-)_n$.
It follows that $(\mu_nY_n)$ belongs to $E_0$, because
$a_-AY=G_--Y_t\in E_0$; similarly for $Z$.
The identity $AY=-Y_{ss}$ holds in distributions by testing against
smooth compactly supported functions and passing in the sine partial
sums. As $Y_t$ and $G_-$ are continuous in space, this second derivative
is continuous. To see the classical conclusion directly, integrate
the continuous right side twice; the difference from $Y$ has zero
distributional second derivative and hence is affine. One elementary
justification of this last assertion is to convolve on every smaller
open interval with a smooth, compactly supported mollifier of integral
one. The convolutions have zero ordinary second derivative and are
affine there; they converge uniformly on compact intervals to the
continuous difference. A uniform limit of affine functions is affine,
as its values at two distinct points determine the coefficients.
Differentiation of this convolution is legitimate because the
mollifier and every derivative have compact support and bounded
integrands. Therefore (3)
is classical. The same reasoning applies to all time derivatives.
Endpoint values are zero by Lemma~\ref{lem:spectral}. The scalar equation gives $b\in C^1$
down to zero, with $b'(0)=0$ since both initial profiles vanish near
the interface. The $E_\alpha$ traces and the rational factors show
that $q_-=Y_s(1)/b$ and $q_+=Z_s(0)/(1-b)$ are analytic in time.
Inverting the coordinate maps proves all claimed local properties.
\end{proof}

\begin{lemma}[Spatial analyticity]\label{lem:spatial}
On every compact $J\subset(0,\delta)$, the local profiles of Lemma~\ref{lem:local}
are analytic in space up to both phase endpoints. Their Taylor series
at the interface represent the actual profiles in an interface
neighborhood.
\end{lemma}
\begin{proof}
Cauchy estimates in $E_\alpha$ and Lemma~\ref{lem:spectral} give, for $q=0,1$,
$\sup_{J\times[0,1]}|\partial_t^n\partial_s^qY|\le C A_0^n n!$.
Equation (3) rearranges as
\[
 Y_{ss}=a(t)Y_t+h(t)sY_s,
 \quad a=b^2/c_-,\quad h=-bb'/c_-. \tag{6}
\]
For $Z$, replace $a$ by $(1-b)^2/c_+$ and $h s$ by
$-(1-b)b'(1-s)/c_+$. On a slightly larger compact interval these
coefficients have bounds $M_aA_0^n n!$, $M_hA_0^n n!$ by Lemma~\ref{lem:local} and \cite[Theorem 3.3.8]{CA}.
Equation (6) and its time derivatives bootstrap continuous spatial
derivatives of every order: at each step its right side is continuous,
and the preceding twice-integration argument applies. Time
differentiation is justified first in distributions, where Lemma~\ref{lem:local} provides
all $E_\alpha$ time derivatives; the resulting continuous
representatives coincide with the classical derivatives.

Choose $A\ge2A_0$. Induction on $q$, simultaneously for all $n$,
gives
\[
 \sup_{J\times[0,1]}|\partial_t^n\partial_s^qY|
 \le C A^n B^q(n+q)! . \tag{7}
\]
For $q\ge2$, use
$\partial_s^{q-2}(sY_s)=s\partial_s^{q-1}Y+(q-2)\partial_s^{q-2}Y$.
In each Leibniz sum, for $m\ge n$,
\[
 \sum_{k=0}^n\binom nk k! (A_0/A)^k(m-k)!
 \le 2m!,
\]
because $\binom nk k!(m-k)!/m!\le1$ and $A_0/A\le1/2$.
The ratios of the three resulting bounds to the right side of (7)
are bounded respectively by
\[
 \frac{2M_aA}{B^2(n+q)},\qquad
 \frac{2M_h}{B(n+q)},\qquad
 \frac{2M_h(q-2)}{B^2(n+q)(n+q-1)}.
\]
Their sum is at most one after increasing $B$, independently of $n,q$.
The right-phase calculation has a minus sign in its last term and
the same absolute bound. Setting $n=0$ in (7), Taylor's remainder
is bounded by $C(B|s-s_0|)^{q+1}$, proving the assertion also at
one-sided endpoints. The remainder formula follows by repeated
integration of the $(q+1)$st derivative along the spatial segment.
\end{proof}

\section{Reconstruction from the interface traces}
For $M,R>0$ and $\sigma\ge1$, write $G(M,R,\sigma)$ for the real
$C^\infty(\mathbb R)$ functions satisfying
\[\sup_t|g^{(n)}(t)|\le M(n!)^\sigma R^{-n}\qquad(n\ge0).\]
A function is bounded Gevrey-$\sigma$ if it belongs to one such class.
For $\sigma=1$ these derivative bounds imply analyticity; for
$\sigma>1$ they allow smooth cutoffs that are flat at their endpoints.

\begin{lemma}[Flatness reconstruction]\label{lem:flatness}
Fix $c>0$, $1\le\sigma<2$ and globally bounded Gevrey-$\sigma$
functions $f,p$ on $\mathbb R$. Set $\beta_0=0$, $\alpha_0=p$,
and recursively
\[
 \beta_{i+1}=c^{-1}(\beta_i'-f\alpha_i),\qquad
 \alpha_{i+1}=c^{-1}(\alpha_i'-f\beta_{i+1}). \tag{8}
\]
Then
\[
 S(t,X)=\sum_{i\ge0}\alpha_i(t)\frac{X^{2i+1}}{(2i+1)!}
       +\sum_{i\ge0}\beta_i(t)\frac{X^{2i}}{(2i)!} \tag{9}
\]
converges on every bounded spatial set with all time and space
derivatives. It is entire in $X$, smooth in $(t,X)$, and satisfies
$S_t=cS_{XX}+fS_X$, $S(t,0)=0$, $S_X(t,0)=p(t)$.
If $f,p$ and all their derivatives vanish at a time, then $S$ vanishes
identically in space at that time.
\end{lemma}
\begin{proof}
Choose a common $R>0$ and positive upper constants such that
$f\in G(M_f,R,\sigma)$, $p\in G(M_\alpha,R,\sigma)$ and
$\beta_0=0\in G(M_\beta,R,\sigma)$. Put
$H=M_\beta/R+M_\alpha M_f$.
With the source parameter $\chi=c^{-1}$,
\cite[Proposition 3.4]{CLT} permits every $\rho>0$ satisfying
\[
 \rho\le \frac{2^{3-\sigma}R}{3},\qquad
 \rho\le \frac{2M_\alpha c}{3H M_f}.
\]
Both upper bounds are strictly positive for any fixed input bounds.
For all integers $k,i\ge0$, the proposition then gives
\[
\begin{aligned}
 \norm{\alpha_i^{(k)}}_\infty
 &\le \frac{M_\alpha c^{-i}}{R^k\rho^i}
 \frac{((k+2i)!)^\sigma}
 {(i!)^{2-\sigma}((2i)!)^{\sigma-1}},\\
 \norm{\beta_{i+1}^{(k)}}_\infty
 &\le \frac{Hc^{-(i+1)}}{R^k\rho^i}
 \frac{((k+2i+1)!)^\sigma}
 {(i!)^{2-\sigma}((2i+1)!)^{\sigma-1}}.
\end{aligned}
\]
For fixed time derivative order, the coefficient estimates of
Proposition 3.4 leave a factor $(i!)^{-(2-\sigma)}$ after division
by the spatial factorial, up to a fixed power of $i$ and a geometric
factor. This proves convergence even for complex $X$ on any disk,
and the same bound applies to every fixed mixed derivative.
\cite[Corollary 3.6]{CLT} permits termwise differentiation. Coefficients of
$X^{2i}$ and $X^{2i+1}$ in $cS_{XX}+fS_X$ are
$c\beta_{i+1}+f\alpha_i=\beta_i'$ and
$c\alpha_{i+1}+f\beta_{i+1}=\alpha_i'$.
This proves the PDE. The interface identities follow from $\beta_0=0$
and $\alpha_0=p$. Each coefficient uses finitely many jets of $f,p$
by (8), proving the final assertion. The cited coefficient estimates impose no input-sign condition.
\end{proof}

\begin{lemma}[Identification]\label{lem:identify}
Suppose on an open time interval in $(0,\delta)$ that $f=b'$ and
$p=q_\pm$, and use (9) with $c=c_\pm$. Then (9), translated by
$X=r-b(t)$, equals the corresponding local temperature throughout
its closed spatial phase.
\end{lemma}
\begin{proof}
For the local temperature expressed in $X$, the equation is
$S_t=cS_{XX}+b'S_X$. Its even and odd spatial derivatives at zero
must therefore obey (8), starting with value zero and slope $q_\pm$.
All mixed derivatives exist by Lemma~\ref{lem:local}--Lemma~\ref{lem:spatial}, so this coefficient calculation
is legitimate. Lemma~\ref{lem:spatial} identifies the local profile with its Taylor series
near zero, while Lemma~\ref{lem:flatness} gives the same Taylor coefficients. Both profiles
are real analytic throughout the connected phase interval, hence
coincide everywhere. For completeness, the latter identity principle
needs no extra PDE theorem: the set of points with a zero neighborhood
for their difference is nonempty and open; at a limit point all
derivatives vanish by continuity, so its convergent Taylor expansion
gives a zero neighborhood. It is thus also closed in a connected
interval. Endpoints follow by continuity.
\end{proof}

\section{The global control}
Products of bounded Gevrey-$\sigma$ functions, affine time rescalings,
and their derivatives are again bounded Gevrey-$\sigma$.
Indeed Leibniz gives
$\binom nk(k!)^\sigma((n-k)!)^\sigma\le(n!)^\sigma$;
the $(n+1)$ terms and geometric factors are absorbed by decreasing
$R$. Differentiation uses $(n+1)^\sigma\le C^n$ after adjusting the
constant for $n=0$. A real analytic function on a neighborhood of a
compact interval has Cauchy bounds $MA^n n!$ there. Multiplying it
by a Gevrey cutoff flat at the interval endpoints, then extending by
zero, gives a globally bounded Gevrey-$\sigma$ function by the same
product bound. These closure properties are used below.

\begin{lemma}[Normalized control construction]\label{lem:global}
For all normalized data above and every $T>0$, there are bounded
Dirichlet controls driving both temperatures to zero and $b$ to
$b_T\in(0,1)$, with a $C^1$ interface strictly inside the rod.
\end{lemma}
\begin{proof}
Let $(\widetilde\theta_\pm,\widetilde b)$ be the solution of Lemma~\ref{lem:local}, with
analytic slopes $\widetilde q_\pm$. Set
$a=\min(\delta,T)/4$ and $t_*=a/2$. Fix $1<\sigma<2$.
By \cite[Lemma 3.2]{CLT}, choose a globally bounded Gevrey-$\sigma$
function $\phi_\sigma:\mathbb R\to[0,1]$, equal to $1$ for $t\le0$
and $0$ for $t\ge1$. All its positive-order derivatives vanish at
both $0$ and $1$. Define
\[
 \chi(t)=\phi_\sigma((t-a)/a),\qquad
 \eta(t)=1-\phi_\sigma((t-a/4)/(a/4)).
\]
Thus $\chi=1$ for $t\le a$, $\chi=0$ for $t\ge2a$,
and $\eta=0$ for $t\le a/4$, $\eta=1$ for $t\ge t_*$.
Define on $[a/4,2a]$, and extend by zero outside,
\[
 f=\eta[\chi\widetilde b'+\chi'(\widetilde b-b_T)],\qquad
 p_-=\eta\chi\widetilde q_-,\qquad p_+=p_--f. \tag{10}
\]
The local quantities are analytic on a neighborhood of this compact
interval because $2a<\delta$. Flatness of the cutoff jets proves
that these zero extensions are globally bounded Gevrey-$\sigma$
functions. In particular there is no regularity claim at the original
time zero for the uncontrolled solution.

For $t\ge t_*$, prescribe
\[
 b(t)=\begin{cases}
 \chi(t)\widetilde b(t)+(1-\chi(t))b_T,&t_*\le t\le2a,\\
 b_T,&t\ge2a.
 \end{cases} \tag{11}
\]
Use $\widetilde b$ for $0\le t\le t_*$.
On $[t_*,a]$ these definitions agree, together with their derivatives.
Both entries in the convex combination (11) lie strictly in $(0,1)$,
and $0\le\chi\le1$; hence so does $b$ at every time.
The flat jets at $2a$ give a smooth joining to the constant path.
For $t\ge t_*$ one has $b'=f$. On the overlap $[t_*,a]$,
$p_\pm=\widetilde q_\pm$, by (1).

Apply Lemma~\ref{lem:flatness} separately with $c_-,p_-$ and $c_+,p_+$, using the same $f$,
and define for $t\ge t_*$
\[
 \theta_-(t,r)=S_-(t,r-b(t)),\quad 0\le r\le b(t),
 \qquad
 \theta_+(t,r)=S_+(t,r-b(t)),\quad b(t)\le r\le1. \tag{12}
\]
Lemma~\ref{lem:flatness} and $b'=f$ give the heat equations after the chain rule.
Their interface values are zero, their slopes are $p_-,p_+$,
and their slope difference equals $f=b'$. Choose the joining time
$t_j=3a/4$, strictly inside $(t_*,a)$, and use the local temperatures
until $t_j$ and (12) afterwards. Lemma~\ref{lem:identify} identifies the two constructions
on the open overlap $(t_*,a)$. The interface defined in (11) also
equals the local interface there. Thus the combined solution satisfies the PDE
across the join with its initial data unchanged.

For $t\ge2a$, all jets of $f,p_-,p_+$ vanish, so both series (12)
are identically zero by Lemma~\ref{lem:flatness}, while $b=b_T$. Since $2a<T$, this gives
the required terminal state. The exterior boundary controls are
the traces $\theta_-(t,0)$ and $\theta_+(t,1)$; they vanish initially
and are continuous later, hence bounded on $[0,T]$.

In the two fixed phase coordinates the combined temperatures are
jointly continuous on compact time--space rectangles and vanish
at the interface. Their zero extensions are therefore jointly
continuous on $[0,T]\times[0,1]$: at a moving-interface crossing
use continuity of the phase-coordinate profiles at their zero
endpoint. The same argument at time zero uses the $E_\alpha$
continuity from Lemma~\ref{lem:local} and the compact initial supports. Uniform
continuity on this rectangle yields continuity into $L^2(0,1)$,
since $\norm{u(t)-u(s)}_2\le\norm{u(t)-u(s)}_\infty$.
Return to the physical variables by
\[
 \ell(t)=Lb(t),\qquad
 y(t,x)=-\frac{\lambda L^2}{d_l}\theta_-(t,x/L),\qquad
 z(t,x)=-\frac{\lambda L^2}{d_r}\theta_+(t,x/L).
\]
The boundary controls are
$v_l(t)=-\lambda L^2\theta_-(t,0)/d_l$ and
$v_r(t)=-\lambda L^2\theta_+(t,1)/d_r$.
For every zero-extended normalized profile $w$,
$\norm{w(\cdot/L)}_{L^2(0,L)}=\sqrt L\,\norm{w}_{L^2(0,1)}$.
These fixed factors preserve boundedness and strong continuity, and
the normalization calculation in (1) restores the physical Stefan law.
This proves Theorem~\ref{thm:control}.
\end{proof}

\paragraph{AI-use statement.}
OpenAI Codex assisted in developing the argument, checking sources and
preparing this manuscript. Independent mathematical review is requested.

\begin{thebibliography}{9}
\raggedright
\bibitem{CLT} B. Colle, J. Loh\'eac and T. Takahashi,
Controllability of the Stefan problem by the flatness approach,
\emph{Systems \& Control Letters} 174 (2023), 105480.
\href{https://doi.org/10.1016/j.sysconle.2023.105480}{doi:10.1016/j.sysconle.2023.105480}.
Source consulted: \href{https://hal.science/hal-03721544v1}{HAL-03721544v1},
submitted July 12, 2022, Lemma 3.2 and Proposition 3.4--Corollary 3.6.
\bibitem{CA} J. Lebl, \emph{Guide to Cultivating Complex Analysis: Working the Complex Field},
version 1.9, July 11, 2026,
\url{https://www.jirka.org/ca/ca.pdf}.
\bibitem{RA} J. Lebl, \emph{Basic Analysis II: Introduction to Real Analysis, Volume II},
version 6.3, May 15, 2026,
\url{https://www.jirka.org/ra/realanal2.pdf}.
\bibitem{AIM258} M. Colbrook, AIM problem collection,
Problem 258, Global constrained null control of a two-phase Stefan interface,
revision c9929805cece044705abaaebe1cebddac1ba11a6.
\href{https://github.com/MColbrook/AIM/blob/c9929805cece044705abaaebe1cebddac1ba11a6/problems/258-global-two-phase-stefan-control.md}{Immutable problem statement}.
\end{thebibliography}
\end{document}
